% Generated by paper/generate_review_sensitivity.py; see review_sensitivity_provenance.json
\begin{table*}[t]
\centering\small
\caption{Post-hoc S1 repository sensitivity for all frozen contrasts.}\label{tab:cluster-sensitivity}
\begin{adjustbox}{max width=\textwidth}\begin{tabular}{lrrrrr}
\toprule
Contrast & Entry mean & Equal-repo mean & Leave-one-repo range & Repo $p$ & Repo $q$ \\
\midrule
Agent--Pack FR & 0.3250 & 0.2228 & [0.2894, 0.3377] & 0.0313 & --- \\
Agent--RAG FR & -0.0250 & 0.0209 & [-0.0463, 0.0049] & 0.6250 & 0.6250 \\
RAG--Pack FR & 0.3500 & 0.2018 & [0.3186, 0.3701] & 0.0625 & 0.1042 \\
Agent--RAG LR & 0.0875 & 0.0533 & [0.0787, 0.0952] & 0.0781 & 0.1042 \\
RAG--Pack LR & 0.2438 & 0.3943 & [0.2245, 0.3238] & 0.0015 & 0.0059 \\
\bottomrule
\end{tabular}\end{adjustbox}
\par\vspace{3pt}\begin{minipage}{\textwidth}\footnotesize Repeats are averaged within entries before grouping the 80 entries into 13 repositories. Exact two-sided sign flips enumerate all $2^{13}=8{,}192$ repository signs using sums of entry differences; this retains the entry-weighted estimand and assumes independently sign-exchangeable repository contributions. The last column is exploratory Benjamini--Hochberg correction across the four secondary comparisons. It does not replace frozen entry-level inference or provide general cluster-robust validity. Equal-repository means target a different weighting.\end{minipage}
\end{table*}
